% Clase 12 --- Qué quiere decir "óhmico", y un ejemplo de lo contrario
% (§4.2 y §6).
% La definición no es "conduce bien" sino una condición geométrica sobre la
% curva: recta Y por el origen. El diodo sirve de contraejemplo porque falla de
% la manera más visible posible --- ni siquiera es simétrico al invertir el
% signo del voltaje.
\begin{center}
\begin{tikzpicture}

  % =============== (a) óhmico ===============
  \begin{axis}[
      fisfig,
      width=7.1cm, height=5.6cm,
      grid=none,
      axis lines=middle,
      xmin=-1.25, xmax=1.25, ymin=-1.25, ymax=1.25,
      xlabel={$\Delta V$}, ylabel={$I$},
      xlabel style={anchor=north west}, ylabel style={anchor=south east},
      xtick=\empty, ytick=\empty,
      axis line style={figgray},
      at={(0,0)},
    ]
    \addplot[figblue, line width=1pt, domain=-1.1:1.1] {0.95*x};
    \node[figlbl, figblue, anchor=west, align=left] at (axis cs:-1.18,0.68)
      {pendiente\\[-0.2em] $1/R$};
    \node[figcarga, fill=figred, minimum size=5pt] at (axis cs:0,0) {};
  \end{axis}
  \node[figlbl, anchor=north, align=center] at (3.0,-1.15)
    {(a) óhmico: una \textbf{resistencia}\\[-0.2em]
     \footnotesize recta \emph{y} por el origen};

  % =============== (b) no óhmico: el diodo ===============
  \begin{axis}[
      fisfig,
      width=7.1cm, height=5.6cm,
      grid=none,
      axis lines=middle,
      xmin=-1.25, xmax=1.25, ymin=-1.25, ymax=1.25,
      xlabel={$\Delta V$}, ylabel={$I$},
      xlabel style={anchor=north west}, ylabel style={anchor=south east},
      xtick=\empty, ytick=\empty,
      axis line style={figgray},
      at={(8.4cm,0)},
    ]
    \addplot[figred, line width=1pt, domain=-1.15:0.95, samples=120]
      {0.04*(exp(2.9*x) - 1)};
    \node[figlbl, figred, anchor=east, align=right, scale=0.9]
      at (axis cs:0.72,0.95) {crece\\[-0.2em] muy rápido};
    \node[figlbl, anchor=north west, align=left, scale=0.9]
      at (axis cs:-1.15,-0.12) {casi no pasa corriente};
  \end{axis}
  \node[figlbl, anchor=north, align=center] at (11.4,-1.15)
    {(b) no óhmico: un \textbf{diodo}\\[-0.2em]
     \footnotesize ni recta ni simétrica};

\end{tikzpicture}\\[0.5em]
{\small\itshape Ser óhmico no es conducir bien: es que la curva sea una recta
\emph{y} pase por el origen. Las dos condiciones importan. Que pase por el
origen dice que sin diferencia de potencial no hay corriente; que sea recta
dice que $R$ es un número, el mismo para toda intensidad ---por eso tiene
sentido hablar de \enquote{la} resistencia de un cuerpo y escribir
$\Delta V = R\,I$---. El diodo incumple lo segundo de forma extrema: su curva
ni siquiera es simétrica, así que la misma diferencia de potencial aplicada al
revés da corrientes que difieren en órdenes de magnitud. Ahí no hay ningún $R$
que reporte: la respuesta depende del punto de trabajo. El diodo es sólo un
ejemplo entre muchas curvas no óhmicas posibles; lo notable es lo contrario,
que tantísimos materiales sí resulten rectos en el rango de campos con que se
trabaja.}
\end{center}
